% !PRINTED 177
\weberpage{177}
Die Auflösung der allgemeinen Gleichungen $2^{\mathrm{ten}}$, $3^{\mathrm{ten}}$ und
$4^{\mathrm{ten}}$ Grades sind also auf Formenprobleme von nur einer Dimension zurückführbar. Wir werden in einem späteren Abschnitte
sehen, dass die allgemeine Gleichung $5^{\mathrm{ten}}$ Grades auf ein binäres
Formenproblem zurückführbar ist, und man kann sich nun als
eine unmittelbare Erweiterung der Aufgabe, die Lösung einer
Gleichung auf reine Gleichungen zurückzuführen, die Frage
stellen: welches ist die geringste Dimensionenzahl eines Formenproblems, durch das sich eine gegebene Gleichung lösen lässt?
Die Aufgabe würde dann so formulirt werden müssen:

\gesperrt{Es sollen aus den Wurzeln einer gegebenen Gleichung rationale Functionen in möglichst kleiner Zahl so gebildet werden, dass sie in homogene lineare Functionen ihrer selbst übergehen, wenn die Wurzeln den Permutationen der Galois'schen Gruppe der gegebenen Gleichung unterworfen werden.}

Auf diese Weise hat F. \gesperrt{Klein} die Aufgabe der algebraischen
Auflösung einer Gleichung erweitert. Er hat, um die Beantwortung der Frage anzubahnen, für die allgemeine Gleichung
$6^{\mathrm{ten}}$ und $7^{\mathrm{ten}}$ Grades bewiesen, dass sie auf quaternäre Formenprobleme zurückgeführt werden können.

Die allgemeine Gleichung $n^{\mathrm{ten}}$ Grades ist, wie wir gesehen
haben, unmittelbar einem Formenproblem von $n$ Dimensionen
äquivalent. Die Frage aber, ob die allgemeine Gleichung $n^{\mathrm{ten}}$
Grades, wenn $n$ grösser als 7 ist, einem Formenprobleme von
weniger als $n$ Dimensionen entspricht, ist noch nicht beantwortet; sie muss wahrscheinlich verneint werden $^{1})$.

\begingroup
\renewcommand{\thefootnote}{1)}
\footnotetext[1]{F. \gesperrt{Klein}, Zur Theorie der allgemeinen Gleichungen $6^{\mathrm{ten}}$ und
$7^{\mathrm{ten}}$ Grades; Mathematische Annalen, Bd. XXVIII, S. 18. Vergl. auch \glqq The
Evanston Colloquium, Lectures on Mathematics by \gesperrt{Felix Klein}\grqq{}, Lecture IX, London und New York, Macmillan and Co. (1894).}
\endgroup

\begin{center}
\S.\ 45.

\gesperrt{Einfluss relativer Invarianten.}
\end{center}

Bei der Definition des Formenproblems im \S.\ 43 haben wir
nur die absoluten Invarianten der Gruppe $S$ benutzt. Nun giebt
es, wie wir gesehen haben, auch Fälle, in denen ausser den ab-
%CONT>>
